\begin{textsample}{POS Dim 1 – ce – Score 52.00 – ce/ce\_ef\_19.txt}  \label{ex:f1_pos_045}
O componente curricular Matemática No \textbf{que} \textbf{diz} respeito à formação pessoal, a Matemática, ao mesmo tempo, demanda e fortalece o pensamento lógico e reflexivo.
O componente Matemática é \textbf{basilar} para a formação de capacidades intelectuais refinadas e, \textbf{certamente}, indispensáveis ao aprendizado das demais áreas do conhecimento. \textbf{Ademais}, a Matemática concilia aspectos do \textbf{raciocínio} indutivo e dedutivo.
De fato, estão combinadas na atividade matemática a busca por generalizações e abstrações, a partir da \textbf{intuição} e de evidências empíricas, e a verificação lógica, em bases dedutivas firmes.
O impulso criativo e intuitivo alia -se, na Matemática, ao rigor lógico e ele representa e envolve aspectos complexos e nobres da formação humanística.
A apresentação da Matemática, em qualquer etapa da educação, não pode ficar circunscrita apenas ao uso cotidiano, social, pragmático, dos conceitos e operações.
Dados do PISA ( 2015 ) \textbf{revelam} \textbf{que} em todas as camadas sociais observam -se dificuldades em relação a operações aritméticas básicas, em sua execução.
Não há explicações \textbf{simples} ou consenso sobre os fatores \textbf{que} levem a esse estado de coisas.
Por vezes, manifesta -se, no ensino de Matemática, a tendência a manipulação formal como \textbf{um} fim em si mesmo, para \textbf{que} cumpramos os ditames curriculares ou o aspecto meramente técnico das competências.
Em algumas ocasiões, evitamos, a qualquer custo a abstração e a generalização, \textbf{que} são características da Matemática, em favor de simplificações \textbf{que} apenas empobrecem a visão da \textbf{disciplina}.
Há estudos \textbf{que} afirmam \textbf{que}, apesar deste componente estar presente nas práticas cotidianas, esta relação tem sido pouco explorada pelo ensino empreendido nas escolas.
Neste espaço, ainda predominam a execução de \textbf{tarefas} repetitivas desconectadas do contexto dos estudantes e das possibilidades de relação com o mundo onde \textbf{vivemos}.
Não devemos confundir abstração com \textbf{mera} habilidade formal do uso da linguagem na forma de \textbf{tarefas} rotineiras.
Tampouco, \textbf{convém} entender \textbf{contextualização} como substituir, no enunciado de \textbf{um} problema, a linguagem formal pela linguagem natural em \textbf{termos} de \textbf{uma} situação “ prática ” ou “ aplicada ”, muitas vezes inverossímil. \textbf{Um} objetivo importante do DCRC é fornecer orientações para a construção dos currículos das escolas.
Nesse sentido, é fundamental entendermos \textbf{que} é difícil superestimar o papel do currículo em indicar as linhas gerais ao longo das quais os conceitos e operações matemáticas devem ser apresentados, seguindo sua disposição lógica intrínseca e apoiando -nos em motivações e \textbf{abordagens} intuitivas cuidadosamente planejadas.
O currículo vai além de \textbf{uma} sucessão cronológica de temas, correspondentes, de algum modo, a conjuntos de competências e habilidades.
Na verdade, deve espelhar as conexões entre os diferentes objetos matemáticos, a interdependência de suas áreas e a estrutura em espiral de complexidade construtiva \textbf{que} caracteriza a Matemática.
Devemos nos assemelhar a \textbf{um} jardim de caminhos \textbf{que} se bifurcam \textbf{que}, partindo dos \textbf{fundamentos}, conduz rapidamente a territórios inexplorados.
O currículo deve ainda promover \textbf{uma} visão interdisciplinar e histórica de aplicações da Matemática, \textbf{que} deponha contra o lugar-comum de \textbf{uma} \textbf{disciplina} difícil, árida, estanque e sem qualquer utilidade óbvia.
Em consonância com o descrito, esses princípios pedagógicos, \textbf{que} expõem o letramento matemático, são definidos em documentos institucionais on-line ( BRASIL, 2017 ) nos seguintes termos : > as competências e habilidades de raciocinar, representar, comunicar e \textbf{argumentar} matematicamente, de modo a favorecer o estabelecimento de conjecturas, a formulação e a resolução de problemas em \textbf{uma} variedade de contextos, utilizando conceitos, procedimentos, fatos e ferramentas matemáticas. ( BRASIL, 2017 ).
Essa definição está alinhada com o conceito estabelecido pela OCDE nos documentos \textbf{que} norteiam o PISA : > Letramento matemático é a capacidade de formular, empregar e interpretar a matemática em \textbf{uma} série de contextos, o \textbf{que} inclui raciocinar matematicamente e utilizar conceitos, procedimentos, fatos e ferramentas matemáticos para descrever, explicar e prever fenômenos.
Isso ajuda os indivíduos a reconhecer o papel \textbf{que} a matemática \textbf{desempenha} no mundo e faz com \textbf{que} cidadãos construtivos, engajados e reflexivos possam fazer julgamentos bem fundamentados e tomar as decisões necessárias ( OCDE, 2012 ).
A condição de letramento é \textbf{detalhada} em \textbf{termos} do desenvolvimento de competências e habilidades, tais como descritas na BNCC em Matemática.
De fato, é o conjunto definido de competências e habilidades na BNCC \textbf{que} estrutura a proposta curricular nacional, ano a ano.
Nesse sentido, tanto as diretrizes nacionais quanto a presente proposta são \textbf{baseadas} em indicações claras do \textbf{que} os alunos devem “ saber ” ( ou seja, competências \textbf{que} envolvem a constituição de conhecimentos, habilidades, atitudes e valores ) e, sobretudo, do \textbf{que} devem “ saber fazer ” ( habilidades \textbf{que} consideram a mobilização das competências para resolver demandas complexas da vida cotidiana, do pleno exercício da cidadania e do mundo do trabalho ).
Dessa forma, o letramento em matemática significa observar o desenvolvimento de diferentes habilidades de relação com o mundo, tais como : ler e compreender informações do mundo presentes em documentos diversos ; analisar e interpretar criticamente dados encontrados nas mais diversas notícias em meios como jornais, revistas e internet ; analisar e decidir a melhor forma de compra de \textbf{um} produto ; participar de atividades \textbf{que} \textbf{exijam} quantificação e operações diferentes cognitivas, dentre tantas outras habilidades.
A exploração dessas habilidades contribui para a ampliação e a exploração de diferentes habilidades em outras áreas de conhecimento, como na Língua Portuguesa, na Geografia, nas Ciências, bem como na Educação Física e nas Artes.
Em alguns casos, \textbf{necessitam} da interpretação de conteúdo, do cálculo de valores, da \textbf{abordagem} em consumo, da compreensão de espaços e de tantas outras ações \textbf{que} mobilizam distintos conhecimentos e formas de relação com o mundo.
Segundo as referências nacionais, descritas na BNCC, a mobilização destes saberes em atividades cotidianas \textbf{necessita} de estar apoiada em princípios universais, como na ética, nos direitos humanos, na justiça social e na sustentabilidade ambiental.
Para \textbf{que} isso aconteça, é \textbf{preciso}, também, garantir o desenvolvimento físico, social, emocional e cultural dos estudantes ( BRASIL, 2017 ).
A formação no Ensino Fundamental, portanto, deve promover o domínio e a capacidade de utilização de conceitos e de recursos da Matemática, a fim de estabelecer adequada relação com o mundo, dentre outras coisas para compreender, formular e resolver problemas, dentro e fora da escola.
Para o ensino de Matemática, ao longo do Ensino Fundamental, propomos, alinhados à BNCC, a exploração de cinco unidades temáticas : Números, Álgebra, Geometria, Grandezas e Medidas e Probabilidade e Estatística.
Nesta estruturação curricular, as mesmas cinco unidades temáticas sugeridas nas BNCC foram mantidas.
Para cada \textbf{uma}, são indicados objetos de conhecimento, aos quais estão relacionadas competências e habilidades específicas, como explicado anteriormente.
É escusado \textbf{dizer} \textbf{que} esta \textbf{divisão} temática é puramente operacional e não pretende, de modo algum, indicar \textbf{que} os assuntos devam ser abordados de modo estanque, em oposição ao aspecto de unidade e coesão da Matemática, em seus conceitos e \textbf{métodos}.
A unidade temática Números pressupõe o desenvolvimento do pensamento numérico, \textbf{que} engloba a noção de número, da contagem, de ideia de quantidade, da escrita numérica e das notações matemáticas.
Também são exploradas noções de aproximação, proporcionalidade, equivalência e ordem.
No entanto, esta unidade não deve ser dissociada das demais unidades temáticas, como se a matemática pudesse ser repartida sem nenhuma relação entre os demais eixos estruturantes.
Nesta defesa, nos primeiros anos, iniciamos com o processo de contagem, apresentado de modo gradativo em conjunto com a ideia de \textbf{um} sistema numérico posicional.
Essa é \textbf{uma} noção \textbf{extremamente} complexa, como sugere o sinuoso percurso histórico ao longo do qual o sistema decimal foi sendo concebido, aceito e \textbf{amplamente} utilizado.
Por isto, apenas, esta já é \textbf{uma} etapa \textbf{que} dispensa \textbf{um} intenso esforço pedagógico.
A seguir, ainda nos anos iniciais do Ensino Fundamental, apresentamos os conceitos dos números naturais aos racionais, enfatizando a completa equivalência de suas representações na forma de frações, razões ou números decimais e, ainda, sua interpretação em \textbf{termos} de resultados da operação de \textbf{divisão}.
Esses fatos óbvios são reforçados à exaustão na proposta, por entendermos \textbf{que}, se não devidamente enfatizados, geram os mal-entendidos e lacunas na formação básica mais prevalentes nos Ensinos Fundamental e Médio.
Conceitualmente, a apresentação aritmético-geométrica dos números racionais ( frações e números decimais, inclusive ), do primeiro ao quinto anos, é calcada no conceito fundamental de proporcionalidade, \textbf{um} dos grandes temas da presente proposta, em todas as unidades temáticas.
No sétimo ano, toda a construção dos números racionais é replicada para incorporar os números negativos, de modo geometricamente natural, por invocar a noção de simetria na reta numérica.
Adição e subtração, multiplicação e \textbf{divisão} são tratadas como operações indissociáveis \textbf{umas} das outras, cujas propriedades e algoritmos devem ser sempre fartamente \textbf{ilustradas} em \textbf{termos} de modelos geométricos.
A noção subjacente de proporcionalidade e sua correlata, linearidade, vão abrindo caminho a questionamentos \textbf{que} indicam a existência de números irracionais.
Prevemos \textbf{um} trabalho prévio de motivação em torno da construção completa dos números reais.
Esse processo de aproximação aos números irracionais e, portanto, ao conjunto completo de reais é realizado de modo gradual, do sétimo ao nono ano.
É \textbf{um} tema de grande complexidade epistemológica e cognitiva e, portanto, requer essa preparação \textbf{conceitual}.
Mesmo na Aritmética, foi adiada a apresentação plena de potência e raízes para o momento \textbf{que} já podemos tratar com propriedade de números irracionais.
Em resumo, ao longo dos anos iniciais e finais, a/o estudante precisará ser capaz de resolver problemas \textbf{que} envolvam as operações básicas com números naturais e racionais, além de entender os significados dessas operações.
Para tanto, devem saber utilizar estratégias próprias e algoritmos ; usar o cálculo mental e saber \textbf{operar} instrumentos como calculadora e computador.
Esta unidade temática ainda prevê, para os anos finais, o estudo de conceitos básicos de economia e finanças, como taxas de juros, inflação e impostos, com o foco na Educação Financeira dos \textbf{aprendizes}.
A unidade temática Álgebra, por sua vez, dispõem -se a desenvolver o pensamento algébrico a partir dos anos iniciais do Ensino Fundamental, \textbf{que} inclui : generalizar padrões ; estabelecer relação entre grandezas ; modelar e resolver problemas aritméticos ; desenvolver habilidades de observação e de interpretação de regularidades a partir de diferentes representações ( tabular, gráfica, simbólica ) ; e abstrair fenômenos matemáticos.
Assim, conforme a BNCC, as ideias matemáticas fundamentais vinculadas a esta unidade são : equivalência, variação, interdependência e proporcionalidade ( BRASIL, 2017 ).
Dessa forma, é \textbf{preciso} propor atividades \textbf{que} contribuam para o entendimento da igualdade, \textbf{enquanto} \textbf{uma} relação \textbf{que} possui características simétricas e transitivas, estabelecendo relações e comparações entre quantidades \textbf{conhecidas} e desconhecidas, como também, tentar expressar alguns significados para \textbf{uma} expressão numérica, para equações e para inequações.
A exploração da Álgebra, nos anos finais, também antecipa o desenvolvimento da capacidade de \textbf{sistematizar}, representar, analisar e resolver problemas por meio da construção de algoritmos⁹.
A exploração destas capacidades pode ser realizada a partir da resolução de problemas modelados pela linguagem algébrica, considerando o conceito de variável e de estrutura lógica operacional, próprios dos algoritmos.
Na presente proposta curricular, esta unidade temática aparece como \textbf{uma} ponte, entre Aritmética e Geometria, por permitir modelar simbolicamente as relações de proporcionalidade e linearidade \textbf{que}, como \textbf{vimos}, são fulcrais no estudo dos números e das formas.
Inicialmente circunscrita, do primeiro ao quinto anos, basicamente à formalização de problemas inversos a Álgebra é expandida a partir do sexto para tornar -se a linguagem \textbf{que} expressará simbolicamente relações entre variáveis numéricas, inicialmente lineares e, a partir do sétimo ano, também não-lineares ( quadráticas, cúbicas, exponenciais, por exemplo ).
Os conteúdos sobre expressões algébricas são naturalmente articulados à necessidade de resolver equações quadráticas e cúbicas \textbf{que}, por sua vez, são a formulação algébrica dos problemas geométricos de construir -se \textbf{um} quadrado ( respectivamente, cubo ) com dada área ( respectivamente, volume ).
Normalmente, \textbf{tópicos} algébricos como esses são tratados nos livros-texto e currículos como expedientes formais gratuitos, sem vinculação ou interesse \textbf{que} não seja o de fatorar intrincadas expressões algébricas.
Tais relações algébricas quadráticas incluem a solução de equações quadráticas por \textbf{mero} completamento de quadrados.
Essa é \textbf{uma} \textbf{abordagem} natural, integrada ao tema das relações não-lineares entre números racionais ou reais.
Ao contrário da prática usual, não colocamos ênfase alguma na necessidade de \textbf{que} os alunos sejam obrigados a ocuparem -se, por \textbf{um} semestre inteiro, em aplicar a chamada Fórmula de Bhaskara a \textbf{um} sem-número de equações quadráticas.
Nesta gradação curricular, a Álgebra abre caminho para a solução de problemas geométricos em \textbf{termos} de números reais.
Por outro \textbf{lado}, as relações lineares e não-lineares entre números reais, bem como as razões trigonométricas no nono ano em \textbf{termos} do conceito algébrico-analítico de funções reais.
A Geometria, por sua vez, envolve o estudo da exploração do espaço ( figuras, formas e relações espaciais ) e de procedimentos necessários para resolver problemas do mundo físico e de diferentes áreas do conhecimento.
Espelhando no \textbf{que} ocorre em Aritmética e Álgebra, os objetos geométricos e suas propriedades são apresentados, inicialmente, com \textbf{forte} apelo intuitivo.
Eles se valem, por exemplo, de modelos concretos, mecânicos e computacionais.
O contato inicial com Geometria convida a aluna, o aluno a pensar não apenas sobre as formas usuais da Geometria Euclidiana e de seus símiles na Natureza e na cultura.
Esse contato também contribui para entender os padrões irregulares \textbf{que} percebemos em ramificações das árvores, em redes neuronais e em outras formas para as quais mesmo \textbf{uma} noção aparentemente \textbf{simples}, como a de dimensão, não se aplica de \textbf{imediato} e sem \textbf{um} árduo processo criativo. ⁹.
Esta capacidade está relacionada ao Pensamento Computacional. “ O Pensamento Computacional compreende abstrações e técnicas necessárias para a descrição e análise de informações ( dados ) e processos.
O domínio destas abstrações e técnicas provê \textbf{uma} habilidade essencial na resolução de problemas, habilidade esta complementar às desenvolvidas em outras áreas como Matemática, física, etc. “ ( SBC, 2018, p. 3 ).
Esta unidade contempla os trabalhos de transformações geométricas e habilidades de construção, representação e interdependência.
São habilidades necessárias ao desenvolvimento de competências relacionadas ao \textbf{raciocínio} e ao pensamento espacial/visual.
Elas são requisitadas, por exemplo, para a leitura de mapa, na interpretação de gráficos estatísticos, nas artes ( pintura, escultura ), na arquitetura, na agricultura e nas engenharias.
A/o estudante desenvolve a competência espacial, quando explora relações de tamanho, \textbf{direção} e posição no espaço ; analisa e compara objetos ; classifica e organiza objetos ; constrói modelos e representações de diferentes situações \textbf{que} envolvem relações espaciais, com desenhos, maquetes, dobraduras e outros.
É importante, também, considerar o aspecto \textbf{funcional} \textbf{que} deve estar presente no estudo da Geometria : as transformações geométricas, sobretudo as simetrias ( BRASIL, 2017 ).
Considerando \textbf{que} as ideias matemáticas fundamentais, associadas a essa temática, são, principalmente a construção, representação e interdependência, podemos propor atividades para \textbf{que} o/a estudante ( com seu corpo e/ou objetos ) vivencie situações de natureza espacial para observar, identificar elementos do universo, perceber propriedades, estabelecer relações e isolar variáveis. \textbf{Ademais}, é esperado \textbf{que} a/o estudante seja capaz de representar figuras planas e/ou \textbf{sólidos} geométricos, seja por malha quadriculada, plano cartesiano, ou ainda, softwares de geometria dinâmica.
Os modelos concretos ou computacionais permitem reunir evidências “ experimentais ” de teoremas fundamentais da Geometria, como os \textbf{que} envolvem a soma dos ângulos internos de \textbf{um} triângulo ou as relações de semelhança entre medidas e ângulos determinados por \textbf{um} feixe de paralelas e transversais, entendido não como \textbf{um} \textbf{mero} fato em meio a \textbf{uma} coleção de proposições, mas como o enunciado formal da própria noção de distância.
A validação experimental deve ser problematizada pela professora/pelo professor na discussão comparativa sobre \textbf{método} indutivo e \textbf{método} dedutivo.
A partir disto, prevemos a introdução à noção de demonstração matemática, base também de construções com régua e compasso.
Esses são dois elementos clássicos e indispensáveis do ensino de Geometria, retomados nessa proposta.
Congruência e semelhança são \textbf{vistas} como noções \textbf{que} definem, dentre todas as possíveis geometria, a Geometria Euclidiana.
Outro elemento original nesta unidade temática é a introdução à Geometria Vetorial no nono ano, ao mesmo tempo, \textbf{aproxima} -se a linguagem algébrica e geométrica do Ensino Fundamental dos conceitos usados na Física, no Ensino Médio e no Cálculo.
A unidade temática Grandezas e Medidas tem \textbf{uma} grande importância social, já \textbf{que} as medidas são usadas para quantificar grandezas do mundo físico, sendo fundamental para a compreensão da realidade.
Esta unidade temática completa, de modo indispensável, todo o percurso curricular em Aritmética, Álgebra e Geometria ao prover modelos concretos de interação com as Ciências e o cotidiano.
Portanto, tem grande relevância em termos utilitários, seja no contexto social ou econômico, seja no contexto científico.
Além disso, é \textbf{uma} unidade \textbf{que} se relaciona naturalmente como outras unidades, inclusive de outros componentes curriculares : Ciências ( densidade, grandezas e escalas do Sistema Solar, energia elétrica, entre outras ) ou Geografia ( coordenadas geográficas, densidade demográfica, escalas de mapas e guias, entre outras ).
Nesse sentido, a BNCC assevera \textbf{que} essa unidade temática contribui ainda para a consolidação e a ampliação da noção de número, da aplicação de noções geométricas e para a construção do pensamento algébrico.
Sendo assim, as/os estudantes do Ensino Fundamental precisam experienciar a resolução de situações-problema \textbf{que} envolvam grandezas de comprimento, área, volume, capacidade, massa, tempo, temperatura e armazenamento de dados computacionais, por exemplo, utilizando, quando necessário, transformações entre unidades de medida padronizadas mais usuais.
Enfatizamos, porém, \textbf{que} o essencial não é simplesmente, o manuseio formal com múltiplos e submúltiplos de unidades de medida, sem \textbf{que} se desenvolva qualquer \textbf{intuição} sobre mudanças de escala.
O essencial é o claro entendimento do próprio processo de medição, especialmente, quando envolve a relação entre grandezas físicas ( densidade, velocidade, relação capacidade/volume ) ou quando \textbf{diz} respeito a grandes mudanças de escala.
É importante \textbf{que} o estudante tenha alguma \textbf{percepção} das imensas diferenças entre, por exemplo, níveis atômicos, celulares, astronômicos e cosmológicos, essenciais para o entendimento da ciência \textbf{contemporânea}.
Dessa forma, alinhados à BNCC ( BRASIL, 2017 ), argumentamos \textbf{que} essa unidade temática contribui ainda para a consolidação e a ampliação da noção de número, da aplicação de noções geométricas e para a construção do pensamento algébrico.
Considerando \textbf{que} as pessoas precisam compreender as informações \textbf{que} estão a sua volta, a unidade temática Probabilidade e estatística propõe o estudo da incerteza, o \textbf{que} pressupõe a necessidade do desenvolvimento da noção de aleatoriedade \textbf{que} deverá possibilitar \textbf{que} as/os estudantes compreendam \textbf{que} nem todo fenômeno é determinístico ; e do tratamento de dados, \textbf{que} envolve o trabalho com a coleta e com a organização de dados de \textbf{uma} pesquisa.
Para o desenvolvimento do \textbf{que} presume esta unidade temática, é \textbf{preciso} incentivar a verbalização dos estudantes em eventos \textbf{que} envolvem o acaso, possibilitando a construção do espaço amostral.
Além disso, é \textbf{preciso} permitir \textbf{que} os estudantes não apenas participem do desenvolvimento de pesquisas, mas de seu planejamento, de modo \textbf{que} possam desenvolver a noção de amostra, de cruzamento de variáveis, de classificação e da definição do gráfico ( CASTRO ; CASTRO-FILHO, 2015 ).
Este tipo de atividade deverá contribuir para a leitura, a interpretação e a construção de gráficos, bem como para a forma de produção de texto escrito para a comunicação de dados.
Destacamos, como pressuposto, a necessidade de integração destas unidades temáticas, considerando, para a aprendizagem da Matemática, a compreensão e a apreensão do significado e de aplicações de objetos matemáticos.
Portanto, \textbf{salientamos} a importância de alvitrar diferentes temas matemáticos e a utilização de recursos didáticos como malhas quadriculadas, ábacos, jogos, livros, vídeos, calculadoras, planilhas eletrônicas e softwares.
Todavia, é \textbf{preciso} propor, para iniciar o processo de formalização matemática, a utilização destes materiais integrados a situações \textbf{que} proporcionem a reflexão e a \textbf{sistematização}.
Neste sentido, buscamos propiciar às alunas/aos alunos \textbf{uma} visão integrada da Matemática a partir do desenvolvimento das relações existentes entre os conceitos e os procedimentos Matemáticos.
Em resumo, o presente documento parte das competências e habilidades definidas na BNCC para \textbf{assimilar} o contexto regional em seus matizes sociais, culturais e educacionais.
Por outro \textbf{lado}, ensaiamos \textbf{uma} maior amplitude e profundidade das competências e habilidades, considerando as premissas expostas anteriormente a respeito do ensino e aprendizado da Matemática e de sua difusão social como parte vital do conhecimento humano.

% matched lemmas: abordagem, ademais, amplamente, aprendiz, aproximar, argumentar, assimilar, basear, basilar, certamente, conceitual, conhecido, contemporâneo, contextualização, convir, desempenhar, detalhar, direção, disciplina, divisão, dizer, enquanto, exigir, extremamente, forte, funcional, fundamento, ilustrar, imediato, intuição, lado, mero, método, necessitar, operar, percepção, preciso, que, raciocínio, revelar, salientar, simples, sistematizar, sistematização, sólido, tarefa, termos, tópico, um, uma, ver, vivar
\end{textsample}
